% sections/morphology.tex

\section{Morphological and Cytogenetic Notes}

\subsection{Phase-Contrast Morphology}

Brightfield imaging at P5 (40\texttimes, phase contrast) showed a uniform adherent
epithelial monolayer with cobblestone packing, consistent with reference images captured
at banking (P1, Vertex Biorepository archive). No spindle-shaped outgrowths, vacuolation,
or granular debris suggestive of a differing cell population were observed.

\begin{center}
\begin{tikzpicture}[scale=1]
  \def\rx{0.42}
  \def\ry{0.30}
  \foreach \x/\y in {0/0, 0.9/0.15, 1.75/-0.1, 0.4/0.75, 1.3/0.85, 2.5/0.4,
                     0.1/-0.8, 1.0/-0.85, 1.9/-0.7, 2.6/-0.3, 3.3/0.1, 3.2/-0.75}{
    \draw[fill=labteal!25, draw=labnavy, thick] (\x,\y) ellipse ({\rx} and {\ry});
    \fill[labnavy] (\x,\y) circle (1.1pt);
  }
  \draw[rounded corners=4pt, draw=labgray, dashed] (-0.6,-1.25) rectangle (3.9,1.25);
  \node[below, labgray, font=\footnotesize] at (1.65,-1.35) {Schematic: representative monolayer field, P5};
\end{tikzpicture}
\end{center}

\subsection{Modal Karyotype}

A rapid metaphase spread count ($n=20$ spreads, Giemsa banding) was performed as a
supplementary check against gross chromosomal contamination or cross-species mix-up.

\begin{table}[h!]
  \centering
  \renewcommand{\arraystretch}{1.25}
  \begin{tabularx}{\textwidth}{@{} X c @{}}
    \toprule
    \rowcolor{labnavy}
    \textcolor{white}{\textbf{Metric}} & \textcolor{white}{\textbf{Observation}} \\
    \midrule
    Modal chromosome number & 58 (hypertriploid range) \\
    Spreads within $\pm 2$ of mode & 17 / 20 (85\%) \\
    Marker chromosomes noted & 2 stable markers, consistent with archived P1 karyotype \\
    Species confirmation & Human (no non-human centromeric banding pattern detected) \\
    \bottomrule
  \end{tabularx}
\end{table}

The modal count and marker chromosomes are consistent with the karyotype recorded at
original banking, providing corroborating (non-definitive) support for the STR-based
identity determination in Section~3.
